\documentclass[conference]{IEEEtran}
\usepackage{amsmath}
\usepackage{amsfonts}
\usepackage{changepage}
\usepackage{amsthm}
\usepackage{amssymb}
\usepackage{tikz}
\usepackage{graphicx}
\usepackage{algorithm}
\usepackage{algorithmic}
\usepackage{graphicx}
\usepackage{caption}
\usepackage{adjustbox}
\usetikzlibrary{positioning, shapes.geometric, arrows.meta}
\usetikzlibrary{arrows.meta, positioning}
\usepackage{float}
\usepackage{graphicx}
\usepackage[colorlinks, citecolor=blue]{hyperref}
\usepackage{balance}
\usepackage{xcolor}
\usepackage{cleveref}
\usepackage{hyperref}
\usepackage{placeins}
\newtheorem{lemma}{Lemma}
\newtheorem{theorem}{Theorem}
\newtheorem{observation}{Observation}
\newtheorem{definition}{Definition}
\usepackage{changepage}
\usepackage[normalem]{ulem}
\title{Data-Aware Task Scheduling for Approximate Computing of Directed Acyclic Graphs (DAGs) in Federated Fog Systems }
\author{}
\date{}


% ------------ ARTICLE ANNOTATIONS / OTHER ------------ %
\newif\ifannote
 \annotetrue  % comment this line to disable annotations

\ifannote
    \newcommand{\anninsert}[2]{{\color{#1}#2}}
    \newcommand{\anncomment}[3]{{\color{#1}[#2: #3]}}
\else
    \newcommand{\anninsert}[2]{#2}
    \newcommand{\anncomment}[3]{}
    % \renewcommand{\sout}[1]{}
\fi

\newcommand{\TA}[1]{\anncomment{red}{TA}{#1}}
\newcommand{\AF}[1]{\anncomment{violet}{AF}{#1}}
\newcommand{\CG}[1]{\anncomment{orange}{CG}{#1}}
\newcommand{\NN}[1]{\anncomment{blue}{NN}{#1}}
\newcommand{\JW}[1]{\anncomment{brown}{JW}{#1}}

\newcommand{\ta}[1]{\anninsert{red}{#1}}
\newcommand{\af}[1]{\anninsert{violet}{#1}}
\newcommand{\cg}[1]{\anninsert{orange}{#1}}
\newcommand{\nn}[1]{\anninsert{blue}{#1}}
\newcommand{\jw}[1]{\anninsert{cyan}{#1}}
% ------------ END ARTICLE ANNOTATIONS / OTHER ------------ %

\begin{document}

\maketitle

\begin{abstract}
Modern natural disaster management increasingly relies on distributed intelligent systems capable of running complex, time-critical analytics close to where hazards occur. Applications such as real-time flood modeling, wildfire spread prediction, structural health monitoring, and evacuation optimization are often expressed as Directed Acyclic Graphs (DAGs) deployed across edge and fog computing platforms embedded in remote, infrastructure-constrained environments. However, during extreme events where cloud connectivity becomes unreliable or unavailable, local fog systems face severe resource inelasticity, hindering their ability to meet strict latency requirements for life-saving computations.

To address this challenge, this work introduces a data-aware and bandwidth-conscious approximate computing framework designed for distributed intelligent sensing and disaster-response analytics. Our approach enables elastic, cooperative processing by (a) federating geographically proximate fog nodes through high-bandwidth/low-latency inter-fog links, and (b) selectively offloading and approximately executing DAG subgraphs across this federated network to ensure real-time decision support with minimal loss in analytical fidelity.

We formulate the resource allocation and task-approximation problem using a Mixed-Integer Linear Programming (MILP) model that jointly considers (i) sensor data size and inter-node bandwidth, (ii) topological and environmental constraints of the distributed computing network, and (iii) per-node computational and memory capabilities. To support large-scale disaster scenarios involving complex multi-sensor pipelines, we further propose meta-heuristic relaxations based on genetic algorithms and simulated annealing, enabling scalable near-optimal solutions.

Experiments performed under realistic disaster-response network conditions demonstrate that the proposed strategies significantly improve situational awareness and processing elasticity, allowing distributed fog systems to meet stringent deadlines while preserving high analytical accuracy. These results highlight the potential of federated fog intelligence as a foundation for resilient, real-time natural disaster management.

Index Terms— Natural Disaster Management, Distributed Intelligence, Federated Fog Systems, Directed Acyclic Graphs, Approximate Computing, Edge Computing, Mixed-Integer Linear Programming, Genetic Algorithms, Simulated Annealing, Real-Time Environmental Monitoring.

\section{Introduction}

The Fourth Industrial Revolution (Industry 4.0) has transformed industrial operations through the integration of automation, artificial intelligence, and real-time data analytics. However, deploying these advancements in remote or extreme environments—such as offshore platforms, deep-sea facilities, or space stations—introduces significant challenges due to high communication latency, intermittent connectivity, and constrained computational resources \cite{wang2021adaptive}. These limitations necessitate decentralized computing paradigms capable of operating autonomously and efficiently without continuous cloud access \cite{dragoni2017microservices}.

Federated fog computing has emerged as a promising solution to enhance computational elasticity by interconnecting multiple fog nodes \cite{mattia2023fogcontrol}. Unlike isolated fog systems, a federated fog environment enables cooperative scheduling, data sharing, and adaptive workload balancing across geographically distributed nodes \cite{provisioning2023}. This architecture is particularly critical for real-time and mission-sensitive applications such as disaster response, predictive maintenance, and AI-driven industrial analytics, where even small delays can lead to system degradation or operational failure \cite{fgcs2024}.

In these environments, workflows are typically modeled as Directed Acyclic Graphs (DAGs), where each node represents a computational task and edges represent data dependencies. However, the execution of such DAG-based workflows is constrained by heterogeneous fog node capabilities and the limited bandwidth of inter-node communication links. To address these challenges, this study introduces a \textbf{data- and bandwidth-aware resource allocation framework} based on Mixed-Integer Linear Programming (MILP), designed to optimally partition and distribute DAG tasks across federated fog nodes while accounting for task-level data transfer sizes, communication delays, and network bandwidth limitations.

To ensure scalability for large and complex workflows, we further develop relaxation meta-heuristics based on genetic algorithms and simulated annealing \cite{kumar2023workload}, enabling near-optimal solutions within practical time bounds. The proposed approach enhances federated fog efficiency by jointly optimizing computational accuracy, data transfer cost, and execution latency. Overall, this research advances the adaptability, resilience, and scalability of Industry 4.0 systems operating in dynamic, resource-constrained, and data-intensive fog environments \cite{fgcs2024}.


\section{Overview and Motivation}

Modern industrial systems are increasingly adopting DAG-based applications powered by AI-driven tasks, deployed on fog computing platforms (mini–data centers) to achieve real-time autonomous control. This transition is particularly crucial in remote locations, such as offshore oil fields \cite{hussain2024resource}, space stations \cite{rehman2018}, and underwater robotic platforms \cite{kabanov2022}, where cloud access is unreliable or nonexistent, and human operators are scarce. DAG-based applications—including disaster simulations, predictive maintenance, and emergency response coordination—must be executed with strict latency and reliability constraints to ensure system safety and operational continuity \cite{wang2021adaptive, rehman2018, cai2019}.

Industrial workloads in such environments demand real-time data processing under harsh operational conditions, where conventional cloud computing models are impractical. However, even when long-distance communication with cloud providers becomes unreliable, nearby fog systems can often maintain stable, high-bandwidth connectivity. For instance, offshore oil rigs located within a limited radius can exchange data over dedicated private fiber or wireless backhaul networks, allowing them to process and share mission-critical information locally even when disconnected from the cloud. This inter-fog connectivity opens an opportunity to form a \textit{federated fog} network that combines local computation and inter-node communication to ensure both resilience and performance in data-driven industrial applications \cite{mattia2023fogcontrol}.

Despite this potential, most fog computing deployments remain isolated and \textit{data-unaware}, focusing solely on computational scheduling without considering the cost of data transfers between interdependent tasks. In practice, DAG-based industrial workflows involve significant intermediate data exchange among tasks---for example, transferring sensor streams, simulation checkpoints, or AI inference results between nodes. These data movements are constrained by the \textit{bandwidth capacity} of the underlying resource network and the \textit{data size} of each edge in the DAG. Ignoring these transfer dependencies often leads to suboptimal scheduling decisions that cause bandwidth congestion, delayed task starts, or violation of real-time deadlines \cite{salehi2016robust}.

To address these challenges, this research introduces a twofold approach to enhance fog-based processing in remote and bandwidth-constrained environments. First, nearby fog systems are seamlessly federated to provide cloud-like elasticity and shared resource utilization. This federation enables dynamic scaling of DAG-based workloads across multiple nodes while leveraging inter-fog communication links of known bandwidth capacities. Second, we introduce a \textbf{data-aware and bandwidth-conscious DAG partitioning framework} that explicitly considers the data size and transfer cost associated with each dependency edge. By jointly optimizing computation and communication, the proposed framework ensures that high-volume data transfers occur along high-bandwidth paths, minimizing both latency and congestion across the federated fog network.

To achieve this optimization, the study employs Mixed-Integer Linear Programming (MILP) formulations that integrate (i) the computational capacity of fog nodes, (ii) their local memory constraints, (iii) the data size of each inter-task edge, and (iv) the available bandwidth and latency characteristics between nodes. The objective is to minimize total workflow completion time while satisfying task precedence and deadline constraints. To enable scalability for large and dynamic workloads, relaxation meta-heuristics such as genetic algorithms and simulated annealing are developed to approximate near-optimal solutions within practical time limits.

In a federated fog environment, variable link bandwidths, intermittent wireless connectivity, and heterogeneous node capacities further complicate scheduling. Without a data-aware mechanism, network saturation can degrade both computational throughput and energy efficiency, ultimately leading to operational failure in mission-critical settings. By incorporating \textbf{data-awareness} into task scheduling, our approach dynamically adapts to data locality and network conditions, ensuring efficient workflow execution and resilient system performance.

Many modern industrial applications exhibit DAG structures where data-dependent tasks—such as deep learning-based anomaly detection, feature extraction, and large-scale simulation modeling—must be executed with minimal processing delays. Real-time sensor data streams impose strict deadlines that must be met to prevent catastrophic failures. For example, in a disaster-prone offshore oil field, a real-time spill monitoring application involves multiple interdependent processes: acquiring sensor data from pipelines, preprocessing and extracting relevant features, running predictive simulations for spill trajectory modeling, and issuing evacuation alerts when environmental hazards are detected. Each of these stages generates intermediate data that must be rapidly transmitted across fog nodes within the available bandwidth constraints.

To illustrate this, Figure~\ref{fig:disaster_response_workflow} presents a DAG-based disaster-response workflow, demonstrating how data-aware and federated fog computing enables low-latency execution of mission-critical tasks. Each task node represents a computational process with associated data dependencies, while interconnecting edges denote data transfers whose sizes and rates influence scheduling decisions. The figure highlights how intelligent partitioning and data-aware scheduling across federated fog nodes ensure efficient resource utilization, meeting both latency and reliability goals even under bandwidth-limited conditions.
(Directed Acyclic Graph)To illustrate this, Figure~\ref{fig:disaster_response_workflow} presents a DAG-based disaster-response workflow, demonstrating how federated fog computing ensures low-latency execution of mission-critical tasks while optimizing resource allocation.


% \FloatBarrier
% \begin{figure}[H]
%     \centering
%     \vspace{-10pt} % Adjust space above the figure
%     \makebox[\linewidth]{\includegraphics[width=1.1\linewidth, keepaspectratio]{disaster-response-application-updated.png}}
%     \vspace{-230pt} % Adjust space below the figure
%     \caption{A DAG-based disaster-response workflow illustrating the microservices involved in the \textit{disaster-response application} (The system utilizes \textit{federated fog computing} to ensure real-time data collection, AI-powered anomaly detection, disaster simulation, decision support, and coordinated emergency response. This approach optimizes resource allocation and minimizes potential damage).}
%     \label{fig:disaster_response_workflow}
% \end{figure}
% \FloatBarrier


\section{Problem Statement and Contributions}

As depicted in Figure~\ref{fig:dag_mapping}, for real-time disaster response systems, a federated fog computing network (\( FG_1 \)--\( FG_5 \)) is established through wireless communication links between distributed fog nodes. Each fog node, represented as \( FG_i \), interconnects via bandwidth-constrained wireless channels, forming a networked infrastructure that enables dynamic, low-latency task execution and data sharing.

In such environments, disaster-response workflows are structured as Directed Acyclic Graphs (DAGs) (\( T_1 \)--\( T_6 \)), where each node represents a computational task, and edges define both task dependencies and associated data transfers~\cite{dragoni2017microservices}. This DAG (\( T_1 \)--\( T_6 \)) must be efficiently mapped onto the fog network (\( FG_1 \)--\( FG_5 \)) to minimize end-to-end latency while respecting bandwidth limitations and communication delays between nodes.

Each fog node exhibits heterogeneity in computational capacity, memory availability, energy consumption, and network bandwidth, introducing variability in execution time and communication latency. These factors complicate real-time task scheduling and increase the risk of deadline violations. To ensure reliable disaster-response performance, the federated fog network must dynamically adapt to resource fluctuations, heterogeneous link capacities, and variable workload intensities while maintaining real-time responsiveness.

For effective execution of disaster-response applications, workflow DAGs must be intelligently partitioned and assigned to fog nodes based on each node’s computational capacity, available bandwidth, and energy profile. Mapping a DAG onto the federated fog infrastructure allows parallel execution of independent tasks, reducing response time and improving overall throughput. Figure~\ref{fig:dag_mapping} illustrates this mapping process, where tasks are scheduled across fog nodes to optimize both computation and communication costs in real time.

However, the heterogeneity of fog nodes and the variability of inter-fog bandwidth make resource allocation a complex optimization problem. Poor DAG partitioning can lead to excessive data transfers across low-bandwidth links, creating bottlenecks that degrade performance and increase latency. Thus, a robust, deadline-aware, and data-aware resource allocation mechanism is essential to optimize task scheduling and communication efficiency across the federated fog environment.

Existing methods, including Mixed-Integer Linear Programming (MILP), have been applied to resource allocation in fog and edge systems. While these models can achieve optimal solutions for small-scale deployments, they often struggle to scale efficiently as the number of tasks and fog nodes increases. Moreover, most prior models focus primarily on computational load balancing while neglecting the critical impact of inter-fog data transfer and bandwidth variability on overall performance~\cite{mattia2023fogcontrol}.

To overcome these challenges, this paper proposes a \textbf{hybrid data- and bandwidth-aware optimization framework} that integrates MILP-based global optimization for task allocation with heuristic-based scheduling for runtime adaptability. The proposed approach simultaneously accounts for computational costs, data transfer sizes, and link bandwidth constraints, ensuring that both computation and communication dependencies are jointly optimized. This hybrid design enables scalable, high-performance execution of real-time DAG workflows across federated fog systems while meeting stringent deadline and reliability requirements.


% \begin{figure}[h]
%     \centering
%     \includegraphics[width=1.0\linewidth]{disaster-response-application-v3.png}
%     \caption{Mapping of disaster-response workflow DAG to the federated fog computing infrastructure.}
%     \label{fig:dag_mapping}
% \end{figure}

\subsection{Contributions}

This paper makes the following key contributions:

\begin{itemize}
    \item We present a novel \textbf{data- and bandwidth-aware DAG-to-Fog mapping strategy} that partitions disaster-response workflows and dynamically allocates tasks across federated fog nodes, accounting for computational heterogeneity, bandwidth constraints, and data transfer dependencies.
    
    \item We introduce a \textbf{mathematical optimization framework} based on Mixed-Integer Linear Programming (MILP) to achieve deadline-aware task allocation, minimizing inter-fog communication overhead while maximizing execution efficiency and resource utilization.
    
    \item We develop an \textbf{integer programming-based scheduling model} that integrates computation and communication costs, ensuring real-time execution while maintaining balanced workload distribution across heterogeneous fog nodes.
    
    \item We propose an \textbf{adaptive heuristic scheduling mechanism} employing meta-heuristic techniques such as genetic algorithms and simulated annealing to dynamically reallocate tasks under varying workloads or node failures, thereby enhancing scalability and fault tolerance.
    
    \item We perform \textbf{comprehensive simulation-based evaluations} using disaster-response DAG workloads to validate the proposed framework, demonstrating significant improvements in execution efficiency, latency reduction, and real-time responsiveness across federated fog networks.
\end{itemize}

Overall, these contributions establish a scalable, data-aware, and deadline-sensitive resource allocation framework that enables efficient and resilient execution of real-time industrial and disaster-response applications in heterogeneous federated fog environments.


\section{Related Work}

Fog computing has garnered substantial research attention over the past decade, primarily due to its ability to bring computation closer to data sources and reduce reliance on remote cloud infrastructure. Research efforts have focused on improving \textit{resource allocation}, \textit{task scheduling}, and \textit{optimization strategies} to support latency-sensitive and data-intensive applications in distributed environments. This section reviews key developments in these domains, emphasizing Directed Acyclic Graph (DAG)-based workflow scheduling in heterogeneous fog systems, which remains a central challenge in achieving low-latency and energy-efficient computation.

\subsection{Fog Resource Allocation and Scheduling}
Early fog computing research primarily addressed static resource allocation, where workloads were assigned based on predefined resource capacities without considering dynamic variations in network conditions or computational load. Studies such as \cite{chiou2022emergency, mattia2023fogcontrol} proposed distributed scheduling models to balance workloads between fog and cloud layers, aiming to reduce latency and energy consumption. However, these models often assumed homogeneous resources and neglected inter-fog communication costs, limiting their scalability in heterogeneous and bandwidth-constrained networks.

Recent work has shifted toward dynamic and adaptive scheduling strategies. For instance, heuristic and meta-heuristic approaches—such as genetic algorithms, particle swarm optimization, and ant colony optimization—have been applied to handle the complexity of multi-objective optimization in fog systems \cite{kumar2023workload}. These methods offer flexibility and scalability but may converge to suboptimal solutions when network bandwidth or task dependencies fluctuate significantly. Consequently, there is an increasing need for hybrid frameworks that combine mathematical rigor with runtime adaptability.

\subsection{DAG-Based Workflow Scheduling in Distributed Environments}
DAG-based workflows are widely used to model complex applications in scientific computing, multimedia analytics, and industrial automation \cite{wang2021adaptive, rehman2018, cai2019}. Each node in a DAG represents a computational task, while directed edges capture data dependencies and communication costs between tasks. Efficient DAG scheduling aims to minimize overall completion time (makespan) while satisfying precedence and deadline constraints.

Several studies have investigated DAG scheduling in cloud and edge contexts. For example, \cite{salehi2016robust} introduced a latency-aware workflow model that prioritizes critical paths to reduce end-to-end delay. Similarly, \cite{dragoni2017microservices} explored microservice-oriented DAG decomposition, enabling modular task deployment across heterogeneous nodes. Despite these advancements, most models remain \textit{computation-centric}, focusing on processor allocation and neglecting data transfer and bandwidth considerations—key factors that directly influence workflow latency in fog systems.

\subsection{Data- and Bandwidth-Aware Scheduling in Fog Networks}
With the growing importance of data-intensive applications, recent research has highlighted the role of \textit{data-awareness} in task scheduling. Data-aware models account for task input/output sizes, transfer latency, and available link bandwidth when mapping tasks to resources \cite{mattia2023fogcontrol}. Works such as \cite{provisioning2023} have shown that incorporating data transfer costs into the optimization process significantly improves overall performance, especially when tasks are distributed across geographically dispersed fog nodes.

However, existing data-aware frameworks often assume stable network bandwidth and static topology. In real-world fog deployments, especially under disaster-response or mobile scenarios, bandwidth availability fluctuates due to link interference, congestion, or environmental disruptions. Current approaches fail to dynamically adapt to these variations, leading to performance degradation or task starvation. This gap underscores the need for a \textbf{bandwidth-conscious and deadline-aware DAG scheduling framework} that can jointly optimize computation and communication costs under uncertainty.

\subsection{Federated Fog and Elastic Resource Sharing}
The concept of \textit{federated fog computing} extends traditional fog architectures by interconnecting multiple fog clusters to share computation and storage resources dynamically \cite{fgcs2024}. Federation introduces new opportunities for workload migration and cooperative scheduling but also amplifies challenges related to synchronization, data transfer, and fault tolerance. While some recent studies have proposed hierarchical or peer-to-peer federations for improved elasticity, few have integrated approximate computing and MILP-based optimization to balance computational accuracy and latency constraints across federated nodes.

\subsection{Summary of Gaps}
In summary, existing research provides valuable foundations in resource allocation, DAG scheduling, and data-aware optimization for fog computing. Nonetheless, three key gaps persist:
\begin{enumerate}
    \item Most models remain computation-centric and do not adequately account for inter-fog bandwidth variability or data transfer costs.
    \item Few frameworks address deadline-aware and approximate task execution in federated fog networks.
    \item Scalability and real-time adaptability remain open challenges for large-scale, heterogeneous fog systems.
\end{enumerate}

To bridge these gaps, this work introduces a \textbf{data- and bandwidth-aware hybrid optimization framework} that combines Mixed-Integer Linear Programming (MILP) with heuristic scheduling to achieve scalable, deadline-sensitive task allocation across federated fog environments.


\subsection{Task Scheduling in Fog Computing}

Task scheduling in fog computing presents unique challenges due to its decentralized architecture, heterogeneous resource distribution, and fluctuating network conditions. Several works have explored strategies to optimize task execution and resource utilization in such dynamic environments.

Lit et al.~\cite{provisioning2023} proposed a deep reinforcement learning (DRL)-based model for adaptive resource provisioning in fog environments. Their work emphasized the importance of dynamic scheduling under variable workloads to reduce latency and improve resource efficiency. However, their approach did not fully address inter-node communication overhead, which becomes a dominant factor in federated fog systems.

Ahmad et al.~\cite{ahmad2014workflow} introduced a data-intensive workflow optimization model that partitions workflows across fog nodes to minimize execution time. Although effective for data-centric applications, their method did not consider fog node heterogeneity, which can significantly influence task performance and energy efficiency.

Oluyisola et al.~\cite{oluyisola2022dag} presented a heuristic-based DAG scheduling algorithm that reduced makespan by dynamically assigning tasks to heterogeneous fog nodes. While their approach achieved improved scheduling efficiency, it neglected network congestion and inter-node communication variability—key challenges in federated fog environments.

Ammar et al.~\cite{ammar2022multicloud} proposed a multi-cloud and fog-aware DAG scheduling framework that optimized execution time and minimized communication costs. Despite achieving lower latency, their model exhibited scalability limitations in large-scale federated fog deployments. The same authors later introduced a machine learning-based adaptive scheduler that leveraged historical execution data for runtime prediction. Although it improved scheduling precision, it introduced computational overhead unsuitable for lightweight fog nodes.

Further research extended DAG scheduling to latency- and energy-aware scenarios. Al-Hajji et al.~\cite{alhajji2022reinforcement} employed reinforcement learning for dynamic task prioritization, improving task allocation efficiency. However, the model’s high computational complexity limited its practicality in large-scale networks. Salehi et al.~\cite{Salehi2016} developed a heterogeneity-aware DAG scheduling algorithm that considered varying node capacities, showing better performance under fluctuating workloads but lacking robust fault-tolerance mechanisms.

Recent studies have explored hybrid approaches combining heuristic and AI-driven methods to balance latency, throughput, and energy efficiency in heterogeneous fog environments. Nevertheless, most existing frameworks assume stable network conditions, overlooking bandwidth fluctuations and inter-fog communication costs, which are critical in real-world federated fog systems.

\subsection{Optimization Techniques}

Optimization techniques in fog computing aim to improve task placement, workload balancing, and energy efficiency while minimizing computational and network overhead. Mixed-Integer Linear Programming (MILP) has been widely applied in resource allocation models. Salehi et al.~\cite{Salehi2016} proposed a MILP-based framework for dynamic resource allocation under high workload conditions, introducing probabilistic task pruning to balance computational cost and completion rates. Although effective, MILP models often suffer from scalability limitations in large deployments.

Perin et al.~\cite{perin2022sustainable} introduced a sustainable edge computing framework that leverages renewable energy sources and predictive scheduling to minimize environmental impact. While energy-efficient, this approach overlooked execution deadlines and data dependencies, making it less applicable to real-time industrial applications.

Kang et al.~\cite{kang2021adaptive} proposed a DRL-based task scheduling model that dynamically adjusted assignments based on workload fluctuations. While DRL methods effectively adapt to changing conditions, they often neglect task dependencies inherent in DAG-based workflows, leading to inefficiencies in data-driven applications. Faticanti et al.~\cite{faticanti2020load} developed a load-balancing framework for heterogeneous fog nodes, improving task distribution but lacking prioritization mechanisms for real-time tasks.

Huang et al.~\cite{huang2022adaptive} combined predictive analytics and reinforcement learning for adaptive DAG scheduling, reducing latency while handling dynamic workloads. However, the computational complexity of their approach made deployment in resource-constrained fog environments challenging. Overall, while optimization-based and learning-based strategies have advanced fog scheduling, they often trade off between solution optimality, scalability, and runtime adaptability.

\subsection{Inter-Node Communication in Federated Fog Systems}

Efficient inter-node communication is vital for maintaining low latency and maximizing resource utilization in federated fog systems. Several models have been proposed to improve reliability, scalability, and fault tolerance across interconnected fog nodes.

Calderon et al.~\cite{calderon2018legacy} proposed hybrid communication models to enhance interoperability between fog and cloud infrastructures, achieving improved reliability and data synchronization. Xie et al.~\cite{xie2023decentralized} presented a decentralized communication model that reduced latency and improved network resilience through distributed coordination. However, their method required high bandwidth availability, limiting its applicability in constrained environments.

Peer-to-peer communication frameworks, such as those proposed by Zhao et al.~\cite{zhao2022peertopeer}, focused on task synchronization and cooperative execution across fog nodes. While efficient, these models suffered from high synchronization overhead under heavy workloads. Zhang et al.~\cite{zhang2022faulttolerant} introduced a fault-tolerant communication model that maintained execution continuity during node failures, though at the cost of additional overhead for state replication.

Han et al.~\cite{wang2023latencyaware} developed a latency-aware communication protocol that prioritized low-latency execution in federated fog environments. Their results demonstrated significant improvements for real-time applications; however, the approach was less effective in highly dynamic or congested networks. Overall, inter-node communication remains a key bottleneck in achieving scalable, fault-tolerant federated fog computing.

\subsection{Summary of Findings}

The reviewed studies collectively highlight the complexity of scheduling, optimization, and communication in fog computing environments. While significant progress has been made, three major challenges persist:
\begin{enumerate}
    \item Limited integration of \textbf{data- and bandwidth-awareness} in existing scheduling and optimization frameworks.
    \item Lack of scalability and adaptability in real-time \textbf{federated fog} systems with heterogeneous nodes.
    \item Insufficient mechanisms for balancing \textbf{computation, communication, and energy efficiency} under dynamic conditions.
\end{enumerate}

These gaps motivate the development of the proposed \textbf{data- and bandwidth-aware hybrid optimization framework}, which jointly optimizes computation and communication costs for scalable, deadline-sensitive execution of DAG-based workflows in federated fog environments.



\section{Problem Definition and Modeling}
\label{model}

\subsection{Problem Definition and Assumptions}

This paper addresses the dependent task allocation problem in a federated fog computing environment, where multiple fog nodes collaboratively execute data-intensive workflows represented as Directed Acyclic Graphs (DAGs). The goal is to minimize overall completion time and inter-fog communication overhead while maximizing computational accuracy under resource and bandwidth constraints.

Let \( T = \{T_1, T_2, \dots, T_n\} \) denote the set of tasks to be executed, and \( F = \{f_1, f_2, \dots, f_m\} \) represent the set of available fog nodes. Each task \( T_i \) must be assigned to exactly one fog node \( f_a \). \af{The binary value \( P_{ia} \in \{0,1\} \) indicates whether task \( T_i \) is allowed be assigned to fog node \( f_a \).}

\paragraph*{Task Execution Model.}  
Task \( T_i \) executed on fog node \( f_a \) requires an execution time \( t_{ia} \geq 0 \). If executed fully, \( T_i \) achieves accuracy \( A_i > 0 \). When executed partially, the achieved accuracy decreases linearly with the fraction of the task executed. Each fog node \( f_a \) has its own computational capacity \( C_a \) (e.g., in FLOPS or MIPS), which influences \( t_{ia} \).

\paragraph*{DAG Dependency Model.}  
The dependencies among tasks are represented as a DAG \( G = (T, E) \), where each directed edge \((T_j, T_i) \in E\) indicates that task \( T_i \) requires the output data of task \( T_j \) before it can start execution. Let \( D_{ji} \) denote the amount of data (in MB\af{\sout{ or GB}}) produced by \( T_j \) and required by \( T_i \).  

If both tasks \( T_j \) and \( T_i \) are executed on the same fog node \( f_a \), the data transfer delay is negligible, and the execution of \( T_i \) can start immediately after \( T_j \) completes. However, suppose they are executed on different fog nodes \( f_b \) and \( f_a \), respectively. In that case, the data must be transmitted over the inter-fog link \((f_b, f_a)\) with available bandwidth \( B_{ba} \) (in \af{MB/s}). \af{Hence, \sout{In this case,}} the communication delay \( c_{ji}^{ba} \) between \( T_j \) and \( T_i \) is computed as:

\[
c_{ji}^{ba} =
\begin{cases}
    \dfrac{D_{ji}}{B_{ba}}, & \text{if } f_b \neq f_a, \\
    0, & \text{if } f_b = f_a.
\end{cases}
\]

\paragraph*{Objective.}  
The objective is to assign tasks to fog nodes such that:
1. The total workflow completion time does not exceed a global deadline \( t_{\max} > 0 \), and  
2. The average task accuracy is maximized \af{\sout{while minimizing communication overhead caused by data transfers between fog nodes}}.

Formally, the optimization problem seeks an \af{assignment matrix \( e_{ia} \) and an execution fraction matrix $l_{ia}$ that balance three competing objectives:  
(i) limiting the overall latency (computation + communication),  
(ii) reducing inter-node data transfer costs, and  
(iii) maximizing aggregate task accuracy.}

\paragraph*{System Assumptions.}  
The system operates under the following assumptions:
\begin{itemize}
    \item Tasks are non-preemptive; once a task starts on a fog node, it must complete on that same node.  
    \item Each task is executed on exactly one fog node.  
    \item Task dependencies are strictly enforced; a task cannot begin execution until all its predecessors have completed and transferred required data.  
    \item Multiple incoming data streams to a task can be received concurrently from different fog nodes.  
    \item Each fog node supports isolated execution (via virtualization or containerization), ensuring that concurrent task execution does not interfere with individual processing times \cite{xu2021task}.  
    \item Network bandwidth \( B_{ba} \) between any two fog nodes is finite and may differ across links due to heterogeneous communication technologies (e.g., Wi-Fi, 5G, or fiber).  
\end{itemize}

By incorporating both data volume \( D_{ji} \) and link bandwidth \( B_{ba} \) as explicit model parameters, the proposed formulation captures the trade-off between computational distribution and data transfer overhead—enabling data- and bandwidth-aware optimization across federated fog nodes.




\begin{itemize}
    \item The overall \textbf{accuracy of a workflow} is determined by the \af{\sout{minimum} average} accuracy among all tasks within that workflow\af{\sout{, as the least accurate task constrains the reliability of the final output}}.
    \item \textit{Future work:} We plan to extend this model by calculating the overall workflow accuracy as a \textbf{convolution of individual task accuracies}, capturing cumulative degradation or propagation of uncertainty across dependent tasks in the DAG. \AF{Move this future work to Section XII.}
\end{itemize}


\subsection{NP-Hardness}

The formulated task allocation problem is \textbf{NP-hard}. This complexity arises from the fact that even simplified versions of the problem—such as the classical \textit{Minimum Multiprocessor Scheduling Problem}—are known to be NP-hard \cite{lenstra1990approximation}. Consequently, obtaining exact optimal solutions is computationally infeasible for large-scale instances. Therefore, practical implementations rely on approximate or heuristic methods to achieve near-optimal task allocations within reasonable time constrain


\begin{table}[t!]
    \centering
    \caption{Summary of Model Parameters and Variables}
    \begin{tabular}{|l|l|}
        \hline
        \textbf{Symbol} & \textbf{Description} \\
        \hline
        \( t_{ia} \) & Time to execute completely task \( T_i \) on fog node \( f_a \) \\ \hline
        \( A_i \) & Accuracy of task \( T_i \) when executed completely
        \af{\sout{; otherwise 0}} \\ \hline
        \( t_{\max} \) & Maximum allowable completion time for tasks \\ \hline
        \( c_{ji}^{ba} \) & Comm.~delay from task \( T_j \) on fog \( f_b \) to task \( T_i \) on fog \( f_a \) \\ \hline
        \( P_{ia} \) & Binary value indicating if task \( T_i \) is allowed on fog \( f_a \) \\ \hline \hline
        \( l_{ia}\) & Fraction of task \( T_i \) executed on fog node \( f_a \) \\ \hline
        \( e_{ia} \) & Binary variable indicating if task \( T_i \) is assigned to fog \( f_a \) \\ \hline
        \( y_i \) & Completion time for task \( T_i \) \\ \hline
        \( M_{ia} \) & Earliest start time of task \( T_i \) on fog node \( f_a \) \\ \hline
    \end{tabular}
    \label{tab:variables}
\end{table}

\section{Mathematical MILP Formulation}
\label{sec:MILP}

We present a data- and bandwidth-aware formulation of the problem as a Mixed-Integer Linear Program (MILP). The model optimizes task scheduling and resource allocation in a federated fog network by considering computational latency, communication bandwidth, and workflow accuracy. 

To summarize, the key parameters and decision variables are listed in Table~\ref{tab:variables} \AF{Add to the table $D_{ji}$ and $B_{ba}$}. These include task execution times, inter-fog bandwidth, data transfer sizes, and accuracy measures critical to the optimization framework.

\paragraph*{Model Parameters}
The main input parameters of the model are:
\begin{itemize}
    \item \( t_{ia} \): \af{Full execution} time of task \(T_i\) on fog node \(f_a\).
    \item \( A_i \): Full execution accuracy of task \(T_i\).
    \item \( t_{\max} \): Workflow deadline.
    \item \( D_{ji} \): Amount of data (in MB\af{\sout{ or GB}}) transferred from task \(T_j\) to task \(T_i\).
    \item \( B_{ba} \): Available bandwidth (in \af{MB/s}) between fog nodes \(f_b\) and \(f_a\).
    \item \( c_{ji}^{ba} = \frac{D_{ji}}{B_{ba}} \): Communication delay between fog nodes \(f_b\) and \(f_a\) when transferring data from \(T_j\) to \(T_i\).
    \item \( P_{ia} \in \{0,1\} \): Feasibility indicator — 1 if \(T_i\) is allowed to run on \(f_a\), 0 otherwise.
\end{itemize}

\noindent
The decision variables are:
\begin{itemize}
    \item \( l_{ia} \in [0,1] \): Fraction of task \(T_i\) executed on fog node \(f_a\).
    \item \( e_{ia} \in \{0,1\} \): Binary variable indicating if \(T_i\) is executed on \(f_a\).
    \item \( y_i \in [0, t_{\max}] \): Completion time of task \(T_i\).
    \item \( M_{ia} \in [0, t_{\max}] \): Earliest possible start time of task \(T_i\) if executed on \(f_a\).
\end{itemize}

The set of in-neighbor tasks of \(T_i\) in the DAG is denoted by \(\text{in}_i\).

\subsection{Basic Formulation}
\label{sec:basic}

\noindent
\textbf{Objective Function:}
\begin{equation}
\label{eq:objective}
    \text{Maximize } \frac{1}{n} \sum_{T_i \in T} A_i \max_{f_a \in F} \{ l_{ia} \}
\end{equation}

The objective maximizes the average accuracy of all tasks while respecting resource and deadline constraints. The model captures both computational performance and communication delay due to data transfer over bandwidth-limited links.

\noindent
\textbf{Constraints:}
\begin{align}
    &e_{ia} \geq l_{ia}, \quad \forall T_i \in T, f_a \in F \label{eq:lvse} \\ 
    &\sum_{f_a \in F} e_{ia} = 1, \quad \forall T_i \in T \label{eq:onefog} \\
    &M_{ia} = \max_{T_j \in \text{in}_i} \left( y_j + \sum_{f_b \in F} e_{jb} \cdot c_{ji}^{ba} \right), \quad \forall T_i \in T, f_a \in F
    \label{eq:resMif-noempty} \\
    &y_i = \sum_{f_a \in F} e_{ia} \left( M_{ia} + t_{ia} l_{ia} \right), \quad \forall T_i \in T \label{eq-36}\\
    &y_i \leq t_{\max}, \quad \forall T_i \in T \label{eq:yitmax} \\
    &e_{ia} \leq P_{ia}, \quad \forall T_i \in T, f_a \in F \label{eq:permitted}
\end{align}

\noindent
In Constraint~\ref{eq:resMif-noempty}, \( c_{ji}^{ba} = D_{ji} / B_{ba} \) explicitly captures data transfer latency between fog nodes, thus introducing both \textbf{data-awareness} and \textbf{bandwidth-awareness} into the scheduling model.

\paragraph*{Interpretation.}
\begin{itemize}
    \item Equation~\ref{eq:objective} maximizes the average accuracy of executed tasks.
    \item Equations~\ref{eq:lvse}–\ref{eq:onefog} ensure that each task is executed on exactly one fog node.
    \item Equation~\ref{eq:resMif-noempty} defines the earliest start time \(M_{ia}\) based on communication delays caused by bandwidth-limited links.
    \item Equation~\ref{eq-36} computes task completion time as a combination of computation time and communication delay.
    \item Equation~\ref{eq:yitmax} enforces deadline constraints, ensuring all tasks finish before \(t_{\max}\).
\end{itemize}

For simplicity, if $\text{in}_i = \emptyset$, we assume $\max(\cdot) = 0$, hence $M_{ia} = 0$.

\subsection{Bandwidth-Aware MILP Refinement}

To explicitly integrate the data and bandwidth parameters into the MILP, we replace \( c_{ji}^{ba} \) with its definition and introduce linear \af{constraints linking transfer time, data size, and bandwidth. These constraints replace Equation~\ref{eq:resMif-noempty}:}
\begin{equation}
\label{eq:bandwidth}
    M_{ia} \geq y_j + \sum_{f_b \in F} e_{jb} \cdot \frac{D_{ji}}{B_{ba}}, \quad \forall T_i \in T, f_a \in F, T_j \in \text{in}_i
\end{equation}
\af{These constraints enforce} that inter-fog data transfers are delayed proportionally to both data size and inverse bandwidth, making the scheduling process \textit{bandwidth-aware}.

\subsection{Objective Interpretation}

The MILP simultaneously maximizes the average workflow accuracy while ensuring that:
\begin{enumerate}
    \item Data transfers across bandwidth-constrained links are \af{\sout{minimized} reduced},
    \item Task completion times respect precedence and communication dependencies, and
    \item The total workflow execution completes within \(t_{\max}\).
\end{enumerate}

This dual optimization of computational accuracy and network efficiency ensures realistic scheduling for bandwidth-limited federated fog systems where large data dependencies exist between DAG tasks.



%     Linear transformations do not change the solution of a problem. In the context of optimization problems like the resource allocation problem in fog computing, transforming constraints using linear techniques, such as replacing non-linear constraints with linear approximations (as done in the linearized formulation), does not alter the feasible region or the optimal solution.
    
%     The basic formulation and the linearized formulation share the same set of feasible solutions because the transformations applied to the constraints are equivalent in terms of the solutions they admit. Specifically, the transformation from the original constraint \( x_i \leq \max_j \{l_{ij}\} \) to \( x_i \leq \sum_f l_{if} \) using binary variables \( P_{if} \) preserves the feasibility and optimality conditions of the original problem.
    
%     Therefore, the linearized formulation accurately represents the constraints of the basic formulation and yields identical optimal solutions. Thus, the theorem holds true.
% \end{proof}





% The task scheduling and allocation problem in  fog federation environments is formulated as follows:

% We define the following parameters:
% \begin{itemize}
%     \item \( t_{if} \geq 0 \): Time required to execute task \( T_i \) on fog node \( f \).
%     \item \( A_i > 0 \): Accuracy of task \( T_i \) when executed completely.
%     \item \( t_{\max} > 0 \): Maximum allowable time to complete all tasks.
%     \item \( c_{ji}^{f'f} \geq 0 \): Communication delay between task \( T_j \) on fog \( f' \) and task \( T_i \) on fog \( f \).
%     \item \( P_{if} \in \{0,1\} \): Binary variable indicating whether task \( T_i \) can be executed on fog node \( f \).
% \end{itemize}

% We define the decision variables:
% \begin{itemize}
%     \item \( l_{if} \in [0,1] \): Fraction of task \( T_i \) computed on fog node \( f \).
%     \item \( e_{if} \in \{0,1\} \): Binary variable indicating if task \( T_i \) is assigned to fog node \( f \).
%     \item \( y_i \in [0, t_{\max}] \): Minimum time to complete task \( T_i \).
%     \item \( M_{if} \in [0, t_{\max}] \): Earliest start time of task \( T_i \) on fog node \( f \).
% \end{itemize}

% The objective is to maximize the minimum task accuracy across all tasks:
% \[
% \text{Maximize } \min_i \left( A_i \max_f \{ l_{if} \} \right)
% \]

% \subsection{Constraints}
% The model is subject to the following constraints:
% \begin{align}
%     &\sum_f e_{if} = 1, \quad \forall i \quad \text{(task assignment to one fog node)} \\
%     &M_{if} = \max_{j \in \text{in}_i} \left( y_j + \sum_{f'} e_{jf'} c_{ji}^{f'f} \right), \quad \forall i,f \quad \text{(task start time)} \\
%     &y_i = \sum_f e_{if} \left( M_{if} + t_{if} l_{if} \right), \quad \forall i \quad \text{(task completion time)} \\
%     &y_i \leq t_{\max}, \quad \forall i \quad \text{(deadline constraint)} \\
%     &e_{if} \leq P_{if}, \quad \forall i,f \quad \text{(permitted task assignment)} 
% \end{align}

\section{Approximation via Randomized Rounding}
\label{sec:randomized}

The MILP formulation described in Section~\ref{sec:MILP} can be solved optimally using commercial or open-source solvers; however, due to the NP-hard nature of the problem, such an approach is computationally feasible only for small to moderate instances. To achieve scalable, near-optimal solutions with polynomial-time complexity, we employ a \textbf{randomized rounding} technique~\cite{raghavan1987randomized}, which provides an approximation framework for the data- and bandwidth-aware task allocation problem in federated fog systems.

\subsection{Linear Relaxation}

The first step in the approximation process involves relaxing the Boolean decision variables \( e_{ia} \in \{0,1\} \).  
Each variable \( e_{ia} \) is replaced with a continuous variable \( r_{ia} \in [0,1] \), representing the probability that task \( T_i \) is assigned to fog node \( f_a \).  
This transforms the MILP into a Linear Programming (LP) problem that can be solved optimally in polynomial time.

The relaxed formulation retains the same objective function and constraints, except for the integrality condition, and now satisfies:
\[
\sum_{f_a \in F} r_{ia} = 1, \quad \forall T_i \in T, \quad \text{and} \quad 0 \leq r_{ia} \leq 1.
\]
The optimal objective value from this LP, denoted \( A_{\text{LP}}^* \), serves as an \textbf{upper bound} on the achievable average accuracy in the original MILP.  
Since communication costs \( c_{ji}^{ba} = D_{ji} / B_{ba} \) are linear in both data size and inverse bandwidth, the relaxed problem remains linear and data- and bandwidth-aware.

\subsection{Randomized Rounding}

The second step converts the fractional solution \( r_{ia} \) into a feasible discrete allocation.  
For each task \( T_i \in T \), we interpret \( r_{ia} \) as the probability that \( T_i \) is assigned to fog node \( f_a \).  
A fog node is then selected randomly according to this probability distribution.  
Formally, for each task \( T_i \), the variable \( e_{ia} \) is set as:
\[
e_{ia} =
\begin{cases}
1, & \text{if fog } f_a \text{ is chosen for } T_i, \\
0, & \text{otherwise.}
\end{cases}
\]

This randomized selection ensures that:
\[
\mathbb{E}[e_{ia}] = r_{ia}, \quad \forall T_i \in T, f_a \in F,
\]
and that every task is assigned to exactly one fog node.  
Importantly, because the LP relaxation already incorporates data transfer delays through \( D_{ji}/B_{ba} \), the randomized rounding process preserves the \textbf{expected communication cost structure} derived from the bandwidth of the resource graph.

\subsection{Approximation Phase}

In the final step, the MILP is resolved using the rounded values \( e_{ia} \) as fixed parameters rather than decision variables.  
This converts the problem into a standard LP where the only remaining variables are continuous (e.g., \( l_{ia}, y_i, M_{ia} \)).  
This step computes the task start and completion times based on the chosen fog assignments while respecting the data transfer delays between nodes.

The resulting solution:
\begin{itemize}
    \item is guaranteed to be feasible for all DAG dependency and bandwidth constraints,
    \item achieves an expected average accuracy close to the LP upper bound \( A_{\text{LP}}^* \),
    \item and maintains polynomial-time solvability.
\end{itemize}

Formally, if \( A_{\text{MILP}}^* \) denotes the optimal average accuracy of the original MILP, then the randomized rounding approach guarantees:
\[
\mathbb{E}[A_{\text{RR}}] \geq (1 - \varepsilon) A_{\text{MILP}}^*, \quad \text{for a small } \varepsilon > 0,
\]
where \( \varepsilon \) depends on the variance introduced by the rounding process.  
Empirically, this approximation achieves near-optimal performance with significantly lower computational complexity.

\subsection{Bandwidth-Aware Approximation Implications}

The inclusion of \( D_{ji} \) (data transfer size) and \( B_{ba} \) (inter-fog bandwidth) in the LP ensures that the randomized rounding step is not purely probabilistic but also \textbf{bandwidth-sensitive}.  
Assignments with high expected communication costs (large \( D_{ji} \) or low \( B_{ba} \)) have smaller \( r_{ia} \) values, reducing their selection probability.  
As a result, tasks are probabilistically biased toward fog nodes with favorable data locality and communication links, enhancing both throughput and latency performance in the federated fog environment.

\subsection{Discussion}

Randomized rounding thus serves as a bridge between optimal MILP-based scheduling and heuristic scheduling methods.  
While it does not guarantee global optimality, it provides a scalable, mathematically grounded approximation that remains responsive to real-world constraints such as bandwidth heterogeneity and data-dependent communication latency.

In the next section, we compare this probabilistic approximation with a baseline greedy DAG scheduling approach that uses local decision-making rather than global optimization.


%\section{Greedy/Naive Approach for DAG Scheduling}
\section{Baseline 1: Greedy DAG Scheduling across Federated Fog}
\label{sec:baseline-greedy}
We design a \emph{bandwidth- and data-aware} greedy scheduler that allocates DAG tasks across federated fog nodes by making locally optimal choices. The heuristic explicitly accounts for (i) task dependencies, (ii) execution times on heterogeneous nodes, and (iii) inter-fog communication delays driven by data sizes and link bandwidth.

\subsubsection{Key Components}
\begin{itemize}
    \item \textbf{Tasks and Dependencies.} Each task \(T_i\) has full-accuracy value \(A_i\) and execution time \(t_{ia}\) on node \(f_a\). Dependencies are edges \((T_j,T_i)\) with data size \(D_{ji}\).
    \item \textbf{Fog Nodes.} Each node \(f_a\) is heterogeneous. Inter-fog links \((f_b,f_a)\) have bandwidth \(B_{ba}\), inducing transfer delay \(c_{ji}^{ba}=D_{ji}/B_{ba}\).
\end{itemize}

\subsubsection{Priority Rule}
To prefer tasks that are both important and time-critical, we use a composite priority key:
\[
\pi_i \;=\; \alpha \cdot A_i \;+\; (1-\alpha)\cdot \text{CP}_i,
\]
where \(\text{CP}_i\) is the (estimated) critical-path length from \(T_i\) to a sink under current \(t_{ia}\) and \(c_{ji}^{ba}\), and \(\alpha\in[0,1]\) trades off accuracy importance vs. schedule urgency (default \(\alpha=0.5\)). Ties are broken by larger \(A_i/t_{i\bullet}\) (accuracy density).

\subsubsection{Algorithm Steps}
\begin{itemize}
    \item \textbf{Initialize.} Set \(y_i\gets\infty\) for all tasks. Compute \(\pi_i\) and maintain a ready set of tasks whose predecessors are finished.
    \item \textbf{Node Selection.} For each ready task \(T_i\), evaluate all permitted nodes \(f_a\) and compute earliest start
    \[
    M_{ia} \;=\; \max_{T_j\in \text{in}_i}\!\left( y_j \;+\; \sum_{f_b\in F} e_{jb}\,c_{ji}^{ba}\right),
    \quad c_{ji}^{ba}=\frac{D_{ji}}{B_{ba}}.
    \]
    \item \textbf{Completion Time.} With a current approximation factor \(\texttt{approx}\in(0,1]\), candidate completion is
    \[
    y_i^{(a)} \;=\; M_{ia} \;+\; t_{ia}\cdot \texttt{approx}.
    \]
    Choose \(f_a\) minimizing \(y_i^{(a)}\) and set \(e_{ia}\gets 1\), \(y_i\gets y_i^{(a)}\).
    \item \textbf{Deadline Adaptation.} If \(\max_i y_i > t_{\max}\), reduce \(\texttt{approx}\leftarrow \texttt{approx}-\delta\) (e.g., \(\delta=5\!\times\!10^{-3}\)) and re-evaluate remaining tasks; otherwise accept schedule.
\end{itemize}

\subsubsection{Pseudocode}
\begin{algorithm}[t]
\caption{Bandwidth/Data-Aware Greedy DAG Scheduler}
\label{alg:greedy_scheduler}
\begin{algorithmic}[1]
\STATE $\texttt{approx}\leftarrow 1$, compute priorities $\pi_i$; set $y_i\leftarrow\infty,\ \forall i$
\REPEAT
  \STATE \textbf{while} $\exists T_i: y_i=\infty$ \textbf{do}
    \STATE pick ready task $T_i$ with maximum $\pi_i$
    \STATE $y^\star\leftarrow\infty$, $a^\star\leftarrow\bot$
    \FOR{each $f_a\in F$ with $P_{ia}=1$}
        \STATE $M_{ia}\leftarrow \max_{T_j\in \text{in}_i}\big(y_j+\sum_{f_b\in F} e_{jb}\cdot \tfrac{D_{ji}}{B_{ba}}\big)$
        \STATE $y_i^{(a)}\leftarrow M_{ia}+t_{ia}\cdot \texttt{approx}$
        \IF{$y_i^{(a)}<y^\star$} 
            \STATE $y^\star\leftarrow y_i^{(a)}$; $a^\star\leftarrow a$ 
        \ENDIF
    \ENDFOR
    \STATE assign $T_i$ to $f_{a^\star}$: $e_{ia^\star}\leftarrow 1$, $e_{ib}\leftarrow 0\ \forall b\neq a^\star$, and $y_i\leftarrow y^\star$
  \STATE \textbf{end while}
  \STATE $Y_{\max}\leftarrow \max_i y_i$
  \IF{$Y_{\max}>t_{\max}$} \STATE $\texttt{approx}\leftarrow \texttt{approx}-0.005$ \ENDIF
\UNTIL{$Y_{\max}\le t_{\max}$}
\end{algorithmic}
\end{algorithm}

\paragraph{Complexity.} With topological processing and a binary heap for priorities, the selection is \(O(n\log n)\); per task we scan up to \(m\) nodes and compute a max over in-neighbors, yielding \(O\!\big(\sum_i (\deg^-(i)+m)\big)=O(|E|+nm)\).

\subsubsection{Why This Greedy is a Strong Baseline}
It is (i) deadline-aware via \(\texttt{approx}\), (ii) bandwidth-aware through \(D_{ji}/B_{ba}\), and (iii) heterogeneity-aware via \(t_{ia}\). While still myopic compared to MILP, it captures the key trade-offs that drive latency in federated fog.

\section{Baseline 2: Genetic-Algorithm Metaheuristic DAG Scheduler across Federated Fog}
\label{sec:baseline-ga}
We employ a bandwidth-/data-aware GA that explores task--node assignments and approximation levels, then evaluates each candidate with exact communication delays \(c_{ji}^{ba}=D_{ji}/B_{ba}\).

\subsubsection{Encoding}
Each individual encodes:
\begin{itemize}
    \item \textbf{Fog assignment vector} \(\mathbf{f}\): entry \(f(i)\in F\) with \(P_{i\,f(i)}=1\).
    \item \textbf{Execution fraction vector} \(\boldsymbol{\ell}\): entry \(\ell_i\in[0,1]\) (maps to \(l_{i\,f(i)}\)).
\end{itemize}

\subsubsection{Fitness (Bandwidth/Data-Aware)}
Given \(\mathbf{f},\boldsymbol{\ell}\), we compute completion times by a forward pass:
\[
M_{i\,f(i)}=\max_{T_j\in \text{in}_i}\!\Big(y_j+\tfrac{D_{ji}}{B_{f(j),\,f(i)}}\Big),\quad
y_i=M_{i\,f(i)}+t_{i\,f(i)}\ell_i.
\]
Let \(Y_{\max}=\max_i y_i\). The fitness to \emph{maximize} is
\[
\Phi=\frac{1}{n}\sum_i A_i \ell_i\;-\;\lambda_1\!\cdot\![Y_{\max}>t_{\max}]\;-\;\lambda_2\!\cdot\!\frac{(Y_{\max}-t_{\max})_+}{t_{\max}},
\]
with large \(\lambda_1\) and moderate \(\lambda_2\); infeasible assignments (e.g., \(P_{i\,f(i)}=0\)) receive \(\Phi=-\infty\).

\subsubsection{Genetic Operators}
\begin{itemize}
    \item \textbf{Selection:} tournament or roulette.
    \item \textbf{Crossover:} one-point on \(\mathbf{f}\) and \(\boldsymbol{\ell}\) (independently).
    \item \textbf{Mutation:} reassign \(f(i)\) within \(\{a:P_{ia}=1\}\); jitter \(\ell_i\leftarrow\text{clip}(\ell_i+\mathcal{N}(0,\sigma^2))\).
    \item \textbf{Repair:} if a parent of \(T_i\) maps to a slow link (small \(B\)) and \(D_{ji}\) is large, bias mutation toward co-locating \(T_i\) with \(T_j\) to reduce \(\tfrac{D_{ji}}{B}\).
\end{itemize}

\subsubsection{Pseudocode}
\begin{algorithm}[t]
\caption{Genetic Algorithm for Bandwidth/Data-Aware Fog Scheduling}
\label{alg:genetic}
\begin{algorithmic}[1]
\STATE initialize population with feasible \((\mathbf{f},\boldsymbol{\ell})\) (respect \(P_{ia}\)), using co-location seeds for large \(D_{ji}\)
\REPEAT
    \STATE evaluate $\Phi$ via forward pass using $c_{ji}^{ba}=D_{ji}/B_{ba}$
    \STATE select parents; apply crossover and mutation; repair infeasibilities
    \STATE elitism: keep top-$k$ individuals
\UNTIL{termination (max iters or no improvement)}
\end{algorithmic}
\end{algorithm}

\paragraph{Discussion.} The GA explores global structures (e.g., grouping high-traffic subgraphs) that the greedy may miss, while preserving feasibility and bandwidth-awareness in each fitness evaluation. In practice, a few dozen generations with modest population sizes provide strong solutions at polynomial cost and complement the LP/MILP and randomized rounding approaches.


\section{Performance Evaluation} 
\label{performance_analysis}

This section evaluates the performance of the proposed task allocation strategies for Directed Acyclic Graphs (DAGs) in federated fog systems. The experiments analyze the effect of three primary parameters—slack factor ($\alpha$), the number of fog nodes ($N$), and DAG size—on scheduling performance. Comparative evaluations are conducted across four approaches: Mixed-Integer Linear Programming (MILP), Randomized Relaxation, Greedy, and Metaheuristic (Genetic Algorithm) scheduling. Performance is measured in terms of average accuracy, scalability, and communication efficiency.

\section{Experimental Setup}

To ensure reproducibility, all experiments were conducted using a custom Python-based simulation framework. The federated fog topology and DAG workloads were synthetically generated and serialized in JSON format, containing attributes such as task execution times, dependencies, deadlines, fog node capacities, and link bandwidths.

The experiments explore how variations in deadline constraints ($t_{max}$), slack factors ($\alpha$), and fog network size ($N$) affect performance under dynamic workload and communication conditions. The simulated environment models a real-time fog federation where distributed nodes cooperate to execute dependent tasks while considering bandwidth constraints and heterogeneous resource availability.

\subsection{Workload Generation}

The workloads represent disaster-response workflows modeled as DAGs of varying scales. Each vertex represents a computational task, and each edge encodes a dependency requiring data transfer between fog nodes.

\begin{itemize}
    \item \textbf{DAG Size:}  
    Three categories of DAGs were generated to evaluate scalability:
    \begin{itemize}
        \item Small-scale: 10 tasks
        \item Medium-scale: 100 tasks
        \item Large-scale: 1000 tasks
    \end{itemize}

    \item \textbf{Task Dependencies:}  
    DAG density was varied from sparse (low inter-task dependencies) to dense (highly interdependent), influencing inter-node communication load.

    \item \textbf{Deadline Constraints:}  
    Deadlines were derived using the critical path length and scaled by slack factor $\alpha$, allowing evaluation under both strict and relaxed scheduling scenarios.

    \item \textbf{Task Heterogeneity:}  
    Tasks were categorized as compute-intensive or latency-sensitive, representing diverse operational profiles typical of edge analytics and emergency response workloads.
\end{itemize}

\subsection{Simulation of the Federated Fog Environment}

A discrete-event simulation framework modeled the cooperative execution of DAGs across interconnected fog nodes. Each fog node was characterized by computational speed, available bandwidth, and link latency. 

\begin{itemize}
    \item \textbf{Network Size:}  
    Simulations were performed for fog networks of varying sizes:
    \[
    N = \{2, 4, 6, 8, 10\}
    \]

    \item \textbf{Resource Heterogeneity:}  
    Each node was assigned a unique processing rate, simulating realistic variations in CPU performance, queue length, and task concurrency.

    \item \textbf{Communication Model:}  
    Link bandwidths ($B_{ba}$) and data transfer sizes ($D_{ji}$) determined communication delays as $c_{ji}^{ba} = D_{ji}/B_{ba}$.  
    This bandwidth-aware formulation ensures that offloaded tasks incur realistic data-transfer delays.

    \item \textbf{Task Offloading Policy:}  
    Tasks were either executed locally or offloaded to remote nodes based on current workload, available bandwidth, and the selected scheduling policy.
\end{itemize}

\subsection{Experimental Design}

The experimental study investigated the following aspects:

\begin{itemize}
    \item \textbf{Effect of Deadline ($t_{max}$):}  
    With $N=10$ fixed, $t_{max}$ was varied to study its influence on task completion, average accuracy, and deadline adherence.

    \item \textbf{Impact of Slack Factor ($\alpha$):}  
    Slack factors $\alpha = \{1.0, 1.2, 1.3, 1.5, 2.0\}$ were used to analyze flexibility versus efficiency trade-offs in meeting deadlines:
    \[
    t_{max} = T_{\text{critical}} + \alpha \cdot T_{\text{slack}}
    \]

    \item \textbf{Scalability with Fog Network Size:}  
    Large DAGs (1000 tasks) were executed with increasing $N$ to assess how network size affects makespan and resource utilization.

    \item \textbf{Max-Accuracy Baseline:}  
    A relaxed configuration ($t_{max} \to \infty$) served as the upper bound for achievable average accuracy across all approaches.
\end{itemize}

\subsection{Baseline Scheduling Methods}

To provide a comparative evaluation, the proposed methods were tested against the following baselines:

\begin{itemize}
    \item \textbf{Fixed Scheduling:}  
    Uses a static slack factor ($\alpha = 1.0$), enforcing strict deadline constraints.
    \item \textbf{Dynamic Slack Allocation:}  
    Adapts $\alpha$ based on workload and network congestion to improve feasibility.
    \item \textbf{Heuristic Scheduling:}  
    Prioritizes tasks by fog node capacity and link latency, approximating a cost-minimization heuristic.
    \item \textbf{Greedy Scheduling:}  
    Allocates each ready task to the least-loaded node, ignoring global dependency structure.
\end{itemize}

\subsection{Performance Metrics}

Experiments were repeated over multiple random DAGs for statistical consistency. The following performance metrics were recorded:

\begin{itemize}
    \item \textbf{Makespan:} Total completion time of all DAG tasks.
    \item \textbf{Resource Utilization:} Ratio of active CPU cycles to total available capacity.
    \item \textbf{Deadline Adherence:} Percentage of tasks completed before $t_{max}$.
    \item \textbf{Communication Overhead:} Aggregate latency and bandwidth cost for inter-node transfers.
    \item \textbf{Scheduling Accuracy:} Average accuracy of all completed tasks.
\end{itemize}

\subsection{Impact of Slack Factor ($\alpha$)}

The slack factor directly scales the execution deadline, affecting latency and resource distribution. The following trends were observed:

\begin{itemize}
    \item Increasing $\alpha$ relaxes deadlines, improving completion rates and average accuracy across all methods.
    \item For $\alpha < 1.0$, strict deadlines cause frequent task preemption and reduced accuracy.
    \item For $\alpha > 2.0$, excessive slack leads to idle resources and increased makespan.
\end{itemize}

Empirical results indicate that a moderate range ($1.2 \le \alpha \le 1.5$) achieves the best trade-off between accuracy and efficiency.  
Figure~\ref{fig:alpha_vs_accuracy} illustrates the relationship between $\alpha$ and average accuracy across scheduling approaches. MILP consistently yields the highest accuracy due to global optimization but scales poorly. Randomized Relaxation achieves near-optimal accuracy with significantly reduced complexity. The Metaheuristic (Genetic Algorithm) approach closely approximates MILP while maintaining good scalability. The Greedy baseline benefits from increased $\alpha$, but its performance plateaus early due to lack of global coordination.



\begin{figure}[h!]
\centering
\vspace{-5mm}
\includegraphics[width=0.5\textwidth]{alpha_vs_accuracy_1.png}
\caption{Impact of Slack Factor ($\alpha$) on Average Accuracy.}
\vspace{-5mm}
\label{fig:alpha_vs_accuracy}
\end{figure}

\subsection{Impact of Fog Nodes}

The number of fog nodes directly affects resource availability for task allocation. Increasing fog nodes reduces contention, leading to higher accuracy. Figure \ref{fig:fog_nodes_vs_accuracy} shows the performance of various methods as the number of fog nodes increases. MILP maximizes the benefits of additional resources, while heuristic methods (Greedy and Relaxation) show diminishing returns. The Relaxation method improves MILP's computational feasibility by providing near-optimal solutions with reduced complexity. The Metaheuristic approach balances accuracy and computational efficiency, leveraging additional resources effectively.

\begin{figure}[h!]
\centering
\vspace{-5mm}
\includegraphics[width=0.5\textwidth]{fog_nodes_vs_accuracy_1.png}
\caption{Impact of Fog Nodes on Average Accuracy. - TODO : 2,4,8,16,32,64}
\vspace{-5mm}
\label{fig:fog_nodes_vs_accuracy}
\end{figure}

\subsection{Impact of Task Graph Size}

As the size of the task graph increases, scalability becomes critical. Figure \ref{fig:task_size_vs_accuracy} illustrates how accuracy trends across methods change with varying DAG sizes. MILP achieves the highest accuracy for small graphs but becomes computationally infeasible for larger graphs. The Relaxation method extends MILP's applicability by simplifying constraints, enabling accurate results for medium-sized graphs. Heuristic methods, while less accurate, scale efficiently and remain computationally viable for large task graphs. The Metaheuristic approach provides a balance, maintaining high accuracy and scalability.

\begin{figure}[h!]
\centering
\vspace{-5mm}
\includegraphics[width=0.5\textwidth]{task_size_vs_accuracy_corrected_1.png}
\caption{Impact of Task Graph Size on Average Accuracy. While MILP theoretically provides maximal accuracy, its performance for larger graphs is limited by computational constraints, leading to reduced accuracy compared to relaxation and metaheuristic approaches.}
\vspace{-5mm}
\label{fig:task_size_vs_accuracy}
\end{figure}

\begin{figure}[h!]
\centering
\vspace{-5mm}
\includegraphics[width=0.5\textwidth]{workflow_failure_vs_accuracy_failure.png}
\caption{Impact of Workflow Node Failures on Average Accuracy. As the percentage of failed nodes increases, the overall accuracy decreases across all methods. MILP maintains higher accuracy under low failure rates but degrades rapidly beyond 30\%. Relaxation and Metaheuristic approaches demonstrate better fault tolerance and sustain accuracy even under high node failure scenarios.}
\vspace{-5mm}
\label{fig:workflow_failure_vs_accuracy}
\end{figure}



\subsection{Scalability Analysis}

Each method demonstrates unique scalability characteristics:
\begin{itemize}
\item \textbf{MILP}: Provides high accuracy for small graphs but experiences reduced accuracy for larger graphs due to computational constraints. Its exponential complexity makes it resource-intensive and less practical for large task graphs.
\item \textbf{Greedy}: Scales efficiently but sacrifices accuracy, particularly for task graphs with complex dependencies. It remains a suitable choice for applications requiring quick, approximate solutions under strict computational constraints.
\item \textbf{Relaxation}: Consistently achieves maximal accuracy across all task graph sizes, demonstrating its scalability and effectiveness. Relaxation serves as a practical alternative to MILP, especially for larger graphs where computational resources are limited.
\item \textbf{Metaheuristic (Genetic Algorithm)}: Provides the best trade-off between scalability and accuracy, handling large graphs effectively. While it achieves near-optimal solutions, it does not match Relaxation’s maximal accuracy but excels in computational efficiency.
\end{itemize}

\subsection{Discussion}

The experimental results highlight key insights about scalability and accuracy:
\begin{itemize}
\item \textbf{MILP}: While theoretically expected to provide maximal accuracy, it experiences reduced performance for larger graphs due to computational constraints. This makes it less practical for real-world large-scale applications.
\item \textbf{Relaxation}: Demonstrates consistent maximal accuracy across all graph sizes. It bridges the gap between MILP’s theoretical guarantees and practical applicability, making it the most reliable method for achieving both accuracy and scalability.
\item \textbf{Greedy and Metaheuristic}: Greedy offers quick, approximate solutions but sacrifices accuracy, while Metaheuristics balance scalability and accuracy well. Metaheuristics are robust and computationally efficient for handling large task graphs, though they do not achieve maximal accuracy.
\end{itemize}

\subsection{Addressing Scalability Challenges}

While Relaxation consistently achieves maximal accuracy, the scalability challenges of other methods necessitate alternative approaches:
\begin{itemize}
\item Relaxation : Provides a reliable and scalable alternative to MILP, maintaining maximal accuracy without the computational overhead of MILP.
\item Metaheuristics : Offer a practical trade-off between accuracy and scalability, making them effective for large-scale systems.
\end{itemize}

These findings emphasize the importance of selecting methods based on the task graph size and computational resource availability. Relaxation emerges as the optimal method for maintaining accuracy and scalability, particularly for larger graphs.

\subsection{Novelty and Comparison with Existing Literature}

The proposed scheduling algorithms are designed specifically for federated fog environments with dynamic task execution constraints. While our algorithms are well-structured and optimized for the given problem, they have not been directly compared with existing literature. This is primarily due to the fact that our study introduces a new problem formulation that has not been extensively explored in previous research.

Unlike traditional scheduling approaches that focus on static resource allocation in fog computing, our work addresses the challenge of adaptive scheduling in federated fog environments with dynamically changing slack factors, heterogeneous fog nodes, and workload fluctuations. Existing works typically optimize for predefined objectives such as energy efficiency, latency reduction, or cost minimization, whereas our approach balances deadline adherence, resource utilization, and execution accuracy in a dynamically evolving system.

Since this problem has not been explicitly addressed in prior research, direct comparisons with existing scheduling frameworks would be non-trivial. However, our experimental design includes multiple baseline methods such as fixed scheduling, heuristic scheduling, and dynamic slack allocation to ensure a fair evaluation of our approach.

Future work could involve adapting existing DAG scheduling techniques and fog-based resource allocation methods to our problem domain, enabling a more comprehensive comparison with state-of-the-art techniques.

\section{Conclusion and Future Work}
 
This paper presents a Mixed Integer Linear Programming (MILP) model for task allocation in  fog federation environments, optimizing DAG workflows while balancing accuracy and resource efficiency. The findings reveal that Relaxation consistently achieves maximal accuracy, making it ideal for medium to large task graphs, particularly where MILP's computational complexity becomes a limitation. In contrast, Metaheuristics excel in scalability and provide a practical trade-off between accuracy and efficiency, especially for large-scale, real-time applications.

These insights emphasize the importance of aligning method selection with system requirements. Relaxation is best suited for accuracy-critical scenarios, while Metaheuristics offer robust solutions for dynamic and resource-constrained environments. Future research will focus on hybrid approaches to further bridge the gap between accuracy and scalability in federated fog systems.

These future directions address pressing scalability and adaptability challenges, providing a pathway for creating more robust, efficient, and versatile  fog federation systems.

\begin{itemize}
\item \textbf{Dynamic Workload Adaptation}: Investigating real-time adaptation mechanisms for fluctuating workloads in large-scale fog networks.

\item \textbf{Graph Compression Techniques}: Developing innovative approaches to simplify DAG structures, reducing node and edge redundancies while preserving critical dependencies. These methods aim to improve computational efficiency without compromising task scheduling accuracy.

\item \textbf{Graph Partitioning Strategies}: Exploring novel partitioning algorithms tailored for fog systems, ensuring balanced distribution of tasks across nodes and minimizing inter-node communication overhead.
\item \textbf{Integration with Emerging Technologies}: Exploring the use of AI and Machine Learning to train datasets for improved task scheduling and resource optimization.
\end{itemize}
These future directions address pressing scalability and adaptability challenges, providing a pathway for creating more robust, efficient, and versatile  fog federation systems in resource constrained environments.
In the future, we will be positioning this system in real-world infrastructure, where fog computing resources will be deployed on-site to handle real-time data processing. These resources will rely on space-sharing scheduling and isolation mechanisms to ensure efficient task execution without interference. By leveraging local processing capabilities, the system will minimize latency and reduce the dependence on unreliable cloud connections, making it suitable for mission-critical applications in environments like offshore oil rigs, smart cities, and industrial IoT systems. This approach is expected to enhance performance and scalability while addressing challenges associated with remote or resource-constrained settings.
%\end{itemize}

\section{Acknowledgments}

We would like to thank the anonymous reviewers for their valuable feedback and insightful suggestions. Their comments have helped improve the quality and clarity of this work. 

This research is supported by the National Science Foundation (NSF).

\bibliographystyle{ieeetr}
\balance
\bibliography{references}

\end{document}

% Objective:
% \begin{equation}
% \label{eq:objective}
%     \text{Maximize } \frac{1}{n} \sum_{T_i \in T} A_i \max_{f_a \in F} \{ l_{ia} \}
% \end{equation}
% Subject to
% \begin{align}
%     &e_{ia} \geq l_{ia}, \quad \forall T_i \in T, f_a \in F \label{eq:lvse-bis} \\ 
%     &\sum_{f_a \in F} e_{ia} = 1, \quad \forall T_i \in T \label{eq:onefog-bis} \\
%     &M_{ia} = \max_{T_j \in \text{in}_i} \left( y_j + \sum_{f_b \in F} e_{jb}  \cdot c_{ji}^{ba} \right),  \forall T_i \in T, f_a \in F %: \text{in}_i \neq \emptyset
%     \label{eq:resMif-noempty-bis} \\
%     &y_i = \sum_{f_a \in F} e_{ia} M_{ia} +  \sum_{f_a \in F} t_{ia} l_{ia} , \quad \forall T_i \in T  \label{eq-36-bis}\\
%     &y_i \leq t_{\max}, \quad \forall T_i \in T \label{eq:yitmax} \\
%     &e_{ia} \leq P_{ia}, \quad \forall T_i \in T, f_a \in F \label{eq:permitted-bis}
% \end{align}

% The change from \Cref{eq-36} to \Cref{eq-36-bis} is trivial.

% Simplifications: 
% Equation \ref{eq-36}:
% $$
% y_i = \sum_f e_{if} \left( M_{if} + t_{if} l_{if} \right) , \quad \forall i
% $$ 
% can become 
% $$
% y_i = \sum_f e_{if} M_{if} + \sum_f t_{if} l_{if} , \quad \forall i
% $$ 
% because $e_{if}=1$ if and only if $l_{if}>0$.

% Then, we can add a variable $z_{if}=e_{if} M_{if}$, so that
% $$
% y_i = \sum_f z_{if} + \sum_f t_{if} l_{if} , \quad \forall i
% $$ 

% Then, $z_{if}=e_{if} M_{if}$ can be replaced by
% \begin{align}
%     z_{if} \leq t_{\max} \cdot e_{if} \\
%     z_{if} \leq M_{if} \\
%     z_{if} \geq M_{if} - (1 - e_{if}) t_{\max}
% \end{align}

% \subsection{Parameters}
% \begin{itemize}
%     \item $t_{if}$: Time to run $T_i$ on fog $f$
%     \item $A_i$: Accuracy if $T_i$ is run completely
%     \item $t_{\max}$: Max completion time
%     \item $c_{ij}^{f'f}$: Communication time between $T_i$ and $T_j$, where $T_i$ is executed in fog $f$ and $T_j$ is executed in fog $f'$
%     \item $P_{if} \in \{0,1\}$ : Binary variable indicating if task $T_i$ is permitted to be executed in fog $f$
% \end{itemize}
% Notes: 
% \begin{itemize}
%     \item $I[P]$ is the indicator function whose value is $0$ if predicate $P=False$ and $1$ if $P=True$.
%     \item $\text{in}_i$ is the set of in-neighbors of $T_i$.
% \end{itemize}

% \subsection{Variables}
% \begin{itemize}
%     \item $l_{if} \in [0,1]$: Fraction of $T_{i}$ computed in fog $f$
%     \item $e_{if} \in \{0,1\}$: Whether $T_{i}$ is executed in fog $f$
%     \item $y_{i} \in [0,t_{\max}]$: Min time to complete $T_i$ across all fogs
%     \item $M_{if} \in [0,t_{\max}]$: Soonest start time of $T_i$ on fog $f$ for each task 
%     \item $z_{if}$: auxilliary variable representing the product $z_{if}=e_{if} M_{if}$ 
% \end{itemize}

% \subsection*{Objective}
% \begin{equation}
%     \text{Maximize } \sum_i A_i \sum_f l_{if}  
% \end{equation}
% \subsection*{Constraints}

% \begin{align}
%     &l_{if} \leq e_{if}, \forall i, f \\ 
%     &\sum_f e_{if} \leq 1, \forall i \\
%     &M_{if} \geq y_j + \sum_{f'} e_{jf'} \cdot c_{ff'}^{ij}, \forall i, f, f', j \in in_i  \\
%     &y_i = \sum_f z_{if} + \sum_f t_{if} l_{if} , \quad \forall i \\
% %    y_i \leq M_{if} + t_{if} \cdot l_{if}, \forall i, f   \\
%     &z_{if} \leq t_{\max} \cdot e_{if}, \forall i,f\\
%     &z_{if} \leq M_{if}, \forall i,f \\
%     &z_{if} \geq M_{if} - (1 - e_{if}) t_{\max}, \forall i,f \\
%     &y_i \leq t_{\max}, \forall i \\
%     &l_{if} \leq P_{if}, \forall i,f 
% \end{align}
